\documentclass[10pt]{beamer}

% Theme phổ biến, có thể thay bằng Warsaw, Madrid, Copenhagen,...
\usetheme{Madrid}
\usecolortheme{default}
\setbeamertemplate{navigation symbols}{}

% Packages cần thiết
\usepackage[vietnam]{babel}
\usepackage[utf8]{inputenc}
\usepackage[T1]{fontenc}
\usepackage{fontspec}
\usepackage{xcolor}
\usepackage{listings}
\usepackage{xspace}

% Cấu hình code Python
\lstset{
  language=Python,
  basicstyle=\ttfamily\small,
  keywordstyle=\color{blue}\bfseries,
  stringstyle=\color{red},
  commentstyle=\color{gray}\itshape,
  numbers=left,
  numberstyle=\tiny\color{gray},
  stepnumber=1,
  numbersep=8pt,
  showstringspaces=false,
  breaklines=true,
  frame=single,
  tabsize=2,
  captionpos=b
}

\title{TỰ ĐỘNG HÓA THU THẬP DỮ LIỆU TÀI CHÍNH \\ (DATA CRAWLING)}
\subtitle{Python for Finance}
\author{Ky Le Hong}
\institute{Hà Nội, Tháng 01/2026}

\begin{document}

% Slide 1 - Title
\begin{frame}
  \titlepage
  \vspace{1em}
  \centering
  \Large Python - trợ lý thu thập data tự động
\end{frame}

% Slide 2 - Vấn đề & Giải pháp
\begin{frame}{Vấn đề \& Giải pháp}
  \begin{columns}[T]
    \column{0.55\textwidth}
      \textbf{Nỗi đau thường gặp:}
      \begin{itemize}
        \item Mở 10 tab web để copy giá mỗi sáng
        \item Dữ liệu bị lệch dòng khi paste vào Excel
        \item Không thể cập nhật realtime khi thị trường biến động
      \end{itemize}

    \column{0.45\textwidth}
      \textbf{Giải pháp Python mang lại:}
      \begin{itemize}
        \item \textcolor{blue}{\textbf{Tốc độ}}: 100 mã chỉ trong $\sim$30 giây
        \item \textcolor{blue}{\textbf{Chính xác}}: Không mỏi, không copy nhầm
        \item \textcolor{blue}{\textbf{Tự động}}: Hẹn giờ chạy trước khi thức dậy
      \end{itemize}
  \end{columns}
\end{frame}

% Slide 3 - Giải phẫu một trang web
\begin{frame}{Giải phẫu một trang web (Web Anatomy)}
  \begin{itemize}
    \item Trình duyệt \quad $\rightarrow$ \quad giống như \textbf{máy in}
    \item Source Code (HTML) \quad $\rightarrow$ \quad giống như \textbf{bộ xương}
    \item Inspect Element (F12) \quad $\rightarrow$ \quad \textbf{máy chụp X-Quang} của Data Analyst
  \end{itemize}

  \vspace{1.2em}
  \centering
  \large\textbf{Screenshot so sánh: Giao diện CafeF $\leftrightarrow$ F12 đầy code}
  % \includegraphics[width=0.9\textwidth]{images/cafef_vs_f12.png}
\end{frame}

% Slide 4 - Chiến thuật: Biến số & Hằng số
\begin{frame}{Chiến thuật: Tìm Biến số \& Hằng số}
  \textbf{Quy luật vàng của Crawling:} Tìm ra những thứ \textbf{lặp lại}

  \vspace{0.8em}
  \begin{block}{Hằng số (Cố định)}
    \begin{itemize}
      \item Tên miền: \texttt{cafef.vn}
      \item Cấu trúc bảng: \texttt{<table class="report">}
    \end{itemize}
  \end{block}

  \begin{block}{Biến số (Thay đổi)}
    \begin{itemize}
      \item Mã cổ phiếu: \texttt{HPG}, \texttt{VIC}, \texttt{VNM}, ...
      \item Số trang: \texttt{page=1}, \texttt{page=2}, ...
    \end{itemize}
  \end{block}

  \vspace{1em}
  \centering
  \Large Ví dụ URL: \texttt{https://cafef.vn/du-lieu/\textcolor{red}{\{BIẾN\_SỐ\}}.chn}
\end{frame}

% Slide 5 - Level 1 Hidden API
\begin{frame}{Level 1 - Con đường bí mật (Hidden API)}
  \begin{description}
    \item[Phương pháp] Tìm API JSON ẩn sau giao diện
    \item[Ưu điểm] Nhanh $\bullet$ Sạch $\bullet$ Dữ liệu dạng Key-Value
    \item[Công cụ] Network Tab (F12) + \texttt{requests}
    \item[Ví dụ điển hình] Lấy lịch sử giá CafeF (.ashx / .json)
  \end{description}

  \vspace{1em}
  \textcolor{teal}{\Large $\Rightarrow$ Không cần phải vật lộn với HTML tag!}
\end{frame}

% Slide 6 - Demo Code API
\begin{frame}[fragile]{Demo Code - Lấy dữ liệu qua API}
\begin{lstlisting}
import requests
import pandas as pd

url = "https://cafef.vn/.../PriceHistory.ashx"
params = {"Symbol": "HPG"}

response = requests.get(url, params=params)
data = response.json()

# Chuyển thẳng thành DataFrame
df = pd.DataFrame(data['Data']['Data'])

# Lưu nhanh
df.to_excel("HPG_price_history.xlsx", index=False)
\end{lstlisting}

\footnotesize\textcolor{gray}{Chỉ cần thay \texttt{Symbol} là lấy được mã khác ngay!}
\end{frame}

% Slide 7 - Level 2 Static HTML
\begin{frame}{Level 2 - Quét Web Tĩnh (Static HTML)}
  \begin{itemize}
    \item Đối tượng: Cophieu68, Vietstock cũ, ...
    \item Đặc điểm: Dữ liệu nằm sẵn trong HTML
    \item Vũ khí tối thượng (1 dòng code!):
  \end{itemize}

  \vspace{1.4em}
  \centering
  \Huge\texttt{pandas.read\_html(url)}

  \vspace{0.6em}
  \Large Tự động tìm tất cả thẻ \texttt{<table>} $\rightarrow$ DataFrame
\end{frame}

% Slide 8 - Thách thức Web Động
\begin{frame}{Thách thức của Web hiện đại}
  \begin{alertblock}{Vấn đề thường gặp}
    \begin{enumerate}
      \item Loading động (React/Vue): Requests chỉ lấy được khung rỗng
      \item Anti-Bot protection: Chặn request không giống người thật
      \item Redirect (Mobile $\rightarrow$ Desktop)
    \end{enumerate}
  \end{alertblock}
\end{frame}

% Slide 9 - Selenium - Robot giả lập
\begin{frame}{Giải pháp: Selenium - Robot điều khiển trình duyệt thật}
  \textbf{Cơ chế hoạt động:}
  \begin{enumerate}
    \item Mở trình duyệt Chrome thật (do code điều khiển)
    \item Truy cập trang web
    \item \textcolor{blue}{\textbf{Ngồi đợi}} (Explicit/Implicit Wait) cho dữ liệu tải xong
    \item Lấy HTML $\rightarrow$ BeautifulSoup / pandas xử lý tiếp
  \end{enumerate}

  \vspace{1.5em}
  \centering
  \Large Robot đang điều khiển Chrome \quad $\rightarrow$ \quad \includegraphics[width=0.4\textwidth]{robot_chrome}
\end{frame}

% Slide 10 - Kỹ thuật Tàng hình
\begin{frame}[fragile]{Kỹ thuật "Tàng hình" (Stealth Mode)}
\begin{lstlisting}[language=Python]
from selenium import webdriver
from selenium.webdriver.chrome.options import Options

options = Options()
options.add_argument("user-agent=Mozilla/5.0 (Windows NT 10.0; Win64; x64) ...")
options.add_argument("--disable-blink-features=AutomationControlled")
options.add_experimental_option("excludeSwitches", ["enable-automation"])
options.add_experimental_option('useAutomationExtension', False)

driver = webdriver.Chrome(options=options)
\end{lstlisting}
\end{frame}

% Slide 11 - Parsing Strategies
\begin{frame}{Kỹ thuật "Phẫu thuật" - Parsing bền vững}
  \begin{block}{Chiến thuật được khuyên dùng}
    \begin{enumerate}
      \item Tìm theo \textbf{ID / Class cố định} (nếu có)
      \item \textbf{Keyword Search} – Tìm theo \textbf{từ khóa} gần dữ liệu
    \end{enumerate}
  \end{block}

  \vspace{1em}
  \centering
  \Large Ví dụ: Tìm chữ \textcolor{red}{"EPS"}, rồi lấy con số \textbf{ngay bên cạnh}
  \vspace{0.5em}
  \large $\Rightarrow$ Bền vững hơn rất nhiều khi web thay đổi giao diện
\end{frame}

% Slide 13 - Quy trình chuẩn
\begin{frame}{Quy trình chuẩn của một Data Engineer}
  \begin{enumerate}
    \item \textbf{Điều tra} (Investigate): Web tĩnh hay động? Có API không? (F12)
    \item \textbf{Chọn chiến thuật}: Requests (dễ) $\rightarrow$ Selenium (khó)
    \item \textbf{Code \& Debug}: Xử lý anti-bot, chờ loading
    \item \textbf{Parsing \& Lưu trữ}: Làm sạch, xuất Excel/CSV/Database
  \end{enumerate}

  \vspace{1.5em}
  \centering
  \Large \textbf{Thứ tự ưu tiên}: API > pandas.read\_html > Selenium
\end{frame}

% Slide 14 - Q&A & Bài tập
\begin{frame}{Q\&A và Bài tập về nhà}
  \begin{block}{Bài tập thực hành}
    Xây dựng bot cập nhật danh mục đầu tư cá nhân \textbf{mỗi sáng}:
    \begin{itemize}
      \item Giá hiện tại
      \item P/E
      \item EPS
    \end{itemize}
    Từ danh sách mã do bạn tự chọn
  \end{block}

  \vspace{1em}
  \begin{block}{Tài nguyên gợi ý}
    \begin{itemize}
      \item Github repo code mẫu
      \item Cheat sheet Pandas + BeautifulSoup + Selenium
    \end{itemize}
  \end{block}
\end{frame}

\end{document}